\documentclass[10pt,a4paper]{amsart}
\usepackage[margin=19mm]{geometry}
\usepackage{amsmath,amssymb,mathtools}
\usepackage[scr=rsfs,cal=euler]{mathalfa}
\usepackage{xfrac}
\usepackage[unicode,hidelinks,bookmarks=false]{hyperref}
\newcommand{\ie}{i.e.\ }
\newcommand{\cf}{cf.\ }
\newcommand{\resp}{respectively\ }
\newcommand{\Subsection}{\subsection}
\providecommand{\nobreakdash}{\nobreak\hbox{-}}
\setlength{\emergencystretch}{2em}
\multlinegap0pt
\usepackage{amssymb}
\usepackage{dsfont}
\usepackage{tcolorbox}
\usepackage{xcolor}
\usepackage{mathtools}
\newtheorem{lemma}{Lemma}
\newtheorem{prop}{Proposition}
\newcommand{\N}{\mathbb{N}}
\newcommand{\ad}{{\rm ad}}
\newcommand{\al}{{\alpha}}
\newcommand{\bt}{{\beta}}
\newcommand{\cF}{{\mathcal F}}
\newcommand{\cG}{{\mathcal G}}
\newcommand{\cK}{{\mathcal K}}
\newcommand{\cM}{{\mathcal M}}
\newcommand{\cR}{{\mathcal R}}
\newcommand{\im}{{\rm i}}
\newcommand{\td}{{\mathtt{d}}}
\title{On the integrability of the Kirchhoff-Pohozaev equation on tori}
\author{Dario Bambusi, Emanuele Haus, Simone Marrocco, Michela Procesi}
\date{Source: arXiv:2609.01891v1 (CC BY 4.0)}
\begin{document}
\maketitle

\noindent\emph{Source and licence.} On the integrability of the Kirchhoff-Pohozaev equation on tori. Version arXiv:2609.01891v1 (CC BY 4.0), distributed by arXiv under the Creative Commons Attribution 4.0 International licence, http://creativecommons.org/licenses/by/4.0/, as recorded in the arXiv licence metadata for this exact version and shown on the abstract page https://arxiv.org/abs/2609.01891v1. Full licence notice: \texttt{licenses/arxiv-2609.01891-CC-BY-4.0.txt}.\par\smallskip

\noindent\textit{Editorial scope.} Counted unit: the article's prop formal-lin (source0.tex lines 1974--1979). The article's complete original proof of it is delivered with it (source0.tex lines 1980--2056, one \texttt{proof} environment of 77 lines). No mathematics is written, repaired or re-derived: the statement and the proof are the article's own lines, in the article's own notation, and no retained line is substituted or re-flowed.  Retained uncounted as the source-local context the proof needs: lines 1902--1906.  Nothing was omitted from any retained span, so the package carries no omission record.  No retained line contains a citation, so the wrapper carries no bibliography and no bibliography entry.  No retained line contains a figure environment, an image-inclusion command or drawing code, so no image is delivered and the package declares no asset dependency.  Editorial formatting only, in addition: an AMS \texttt{amsart} preamble that restates the article's own declarations and macros used by the retained lines, namely the environments \texttt{lemma}, \texttt{prop} on separate counters, together with the standard packages those lines need.  The retained environments print in this document as Lemma 1, Proposition 1 in order of appearance, whereas the article prints the same items with the numbering of its own full document.  The complete native source of the article as distributed in the arXiv e-print for this exact version is bundled unchanged as \texttt{sources/209151/source0.tex}.
\begin{lemma} \label{compo} Given a sequence of generating functions $G_i\in \cF^{{\geq} \td_i}$ with $\td_{i+1} > \td_i\geq 1$ then there exists $\cG\in \cF^{\geq d_1}$ such that the composition
	\[
	\prod_i	e^{\{G_i,\}} = e^{\{\cG,\}}
	\] 
\end{lemma}

\begin{prop}\label{formal-lin}Consider a Family of Hamiltonians satisfying H1 and H2.
	There exists $S\in  \cF^{\geq 1}$ such that for all $k\in \N_{\geq1}$ one has 
	\[
	e^{\{S,\cdot\}} H_k = L_k  + Z_k\,,\quad Z_k \in \cK^{\geq 2}\,.
	\]
\end{prop}

\begin{proof}
	We proceed by induction. Let us assume that, for some $n\geq 1$ there exists an $S_n$ such that
	for all $k$
	\[
	e^{\{S_n,\cdot\}} H_k=:H_{k,n}= L_k  + Z_{k,n}+ R_{k,n}\,,\quad Z_{k,n} \in \cK^{2\leq  \td < n}\,, \quad R_{k,n} \in \cR^{\geq n}\,.
	\]
	In order to make a normal form step we look for a formal change of variables $e^{\{F_n,\cdot\}}$, with $F_n\in \cR^{(n)}$ so that
	\[
	e^{\{F_n,\cdot\}} H_k = L_k  + Z_{k,n+1}+ R_{k,n+1 },,\quad Z_{k,n} \in \cK^{2\leq  \td < n+1}\,, \quad R_{k,n} \in \cR^{\geq n+1}
	\]
	By construction 
	\begin{align*}
		e^{\{F_n,\cdot\} }H_{k,n} &=L_k+ Z_{k,n} + R_{k,n} + \{F_n, L_k\} + \sum_{h=2}^\infty \frac{\ad_{F_n}^{h-1}}{h!} \{F_n,L_k \}
		+ \sum_{k=1}^\infty \frac{\ad_{F_n}^k}{k!}(Z_{k,n} + R_{k,n}) \\
		&= L_k+ Z_{k,n} + \Pi^\cK R_{k,n}  - \sum_{k=1}^\infty \frac{\ad_{F_n}^{k}}{(k+1)!}\Pi^{\cR}  P_{i}  + \sum_{k=1}^\infty \frac{\ad_{F_n}^k}{k!}(Z_{k,n} + R_{k,n})\,.
	\end{align*}
	So we may set
	\[
	Z_{k,n+1}:= Z_{k,n}+%\Pi^{\leq 2 i +2}
	\Pi^{(n)}\Pi^\cK R_{k,n} \,,\quad R_{k,n+1}= e^{\{F_n,\cdot\} }H_{k,n}- L_k - Z_{k,n+1}\,.
	\]
	In order to iterate our inductive step we need to show that we may fix $F_n$ so that for all $k$ 
	\[
	\Pi^{\td\leq  n} R_{k,n+1}= \{F_n,L_k\} + \Pi^\cR R_{k,n}^{(n)} =0\,.
	\]
	Passing to the Taylor coefficients this reads
	\begin{equation}
		\label{homotutte}
		\im \omega_k \cdot (\al-\bt)(F_n)_{\al,\bt} =(R_{k,n})_{\al,\bt} \,,\quad \forall \al\ne \bt\,:\quad |\al|+|\bt| =n+2\,.
	\end{equation}
	Finally the condition $\{H_k,K_j\}=0$ ensures that the equation above is trivial unless
	\begin{equation}
		\label{condk}
		{\al_j-\bt_j = \al_{-j} -\bt_{-j}}
	\end{equation}
	We claim that a  common solution to  the above infinite list of equations is given by
	\begin{equation}
		\label{questo}
		(F_n)_{\al,\bt} =\begin{cases}
			\im \dfrac{(R_{1,n})_{\al,\bt}}{\omega_1 \cdot (\al-\bt)}\,,\qquad \mbox{if}\;  \omega_1 \cdot (\al-\bt)\neq 0\\
			\im 	\dfrac{(R_{2,n})_{\al,\bt}}{\omega_2 \cdot (\al-\bt)}\,,\qquad \mbox{if}\;  \omega_1 \cdot (\al-\bt)=0 \;\mbox{and}\;   \omega_2 \cdot (\al-\bt)\neq 0\\
			\vdots
			\\
			\im 	\dfrac{(R_{k,n})_{\al,\bt}}{\omega_k \cdot (\al-\bt)}\,,\qquad \mbox{if}\;  \omega_1j\cdot (\al-\bt)=0 \; \forall j< k \;\mbox{and}\;   \omega_k \cdot (\al-\bt)\neq 0\\
			\vdots
		\end{cases}
	\end{equation}
	for all $(\al,\bt) \in \cM$ satisfyting \eqref{condk} and  $(F_n)_{\al,\bt}=0$ otherwise.
	To prove our claim, we use the fact that, by the involution hypothesis H1., we have 
	\[
	0= e^{\{S_n,\cdot \}} \{H_{k_1},H_{k_2}\}= \{H_{k_1,n},H_{k_2,n}\}=\{ L_{k_1}  + Z_{k_1,n}+ R_{k_1,n},L_{k_2}  + Z_{k_2,n}+ R_{k_2,n}\}
	\]
	so that, since $\{Z_{k_1,n},Z_{k_2,n}\}=0$ and $R_{k_i,n} \in \cF^{\geq n}$, one has 
	\[
	\Pi^{(n)} \{H_{k_1,n},H_{k_2,n}\} =\{L_{k_1}, R_{k_2,n}^{(n)}\} + \{R_{k_1,n}^{(n)}, L_{k_2}\}=0\,,
	\]
	passing to the Taylor coefficients we have
	\begin{equation}
		\label{commuto}
		\omega_{k_2} \cdot (\al-\bt) (R_{k_1,n})_{\al,\bt} = \omega_{k_1} \cdot (\al-\bt) (R_{k_2,n})_{\al,\bt} \,,\qquad \forall \; |\al|+|\bt| = n+2\,.
	\end{equation}
	Denoting $\ell= \al-\bt$, we recall that, by the independence assumption H2., there exists ${\mathbf k}_0 = {\mathbf k}_0(\al-\bt)$ such that
	\[
	\omega_{h} \cdot (\al-\bt) =0\quad  \forall h< {\mathbf k}_0\,,\qquad  \omega_{{\mathbf k}_0}\cdot (\al-\bt) \neq 0\,.
	\]
	By \eqref{commuto}, for all $k$ and all $ |\al|+|\bt| = n+2$ we have
	\[
	(R_{k,n})_{\al,\bt} =% \big( \omega_{{\mathbf k}_0}\cdot (\al-\bt)\big)^{-1}
	\frac{\omega_{k} \cdot (\al-\bt)}{\omega_{{\mathbf k}_0}\cdot (\al-\bt)} (R_{{\mathbf k}_0,n})_{\al,\bt} \,,\qquad \forall \; |\al|+|\bt| = n+2\,,
	\]
	which shows that $F_n$ defined in \eqref{questo} satisfies \eqref{homotutte}.
	Finally by Lemma \ref{compo} we can define $S_{n+1}\in \cF^{\geq 1}$ so that
	\[
	e^{\{S_{n+1},\cdot\}}= e^{\{F,\cdot \} }e^{\{S_n,\cdot\}}\,,
	\]
	passing to the limit we obtain the desired result.
\end{proof}

\end{document}
